\documentclass[aspectratio=169,11pt]{beamer}

\usetheme{Madrid}
\usecolortheme{default}
\usefonttheme{professionalfonts}
\setbeamertemplate{navigation symbols}{}
\setbeamertemplate{footline}[frame number]

\usepackage{amsmath, amssymb, booktabs}

\title{Technology, Tasks, And Labor-Market Adjustment}
\subtitle{Core Frameworks}
\author{Special Topic 5: Technology, Innovation, and Labor -- Week 1}
\date{}

\begin{document}

\begin{frame}
  \titlepage
\end{frame}

\begin{frame}{Central question}
\begin{itemize}
    \item How should labor economists map technologies into tasks, skills, labor demand, worker adjustment, and equilibrium outcomes?
    \item Technology is not the outcome. Wages, employment, mobility, organization, inequality, and welfare are the labor outcomes.
    \item Week 1 discipline: define the technology object before interpreting a technology coefficient.
\end{itemize}
\end{frame}

\begin{frame}{Why this is labor economics}
\begin{itemize}
    \item Technology changes what firms want workers to do.
    \item Innovation changes which skills are valuable and which tasks can be performed by labor, capital, software, or teams.
    \item Adoption changes hiring, training, supervision, workflow, and job design inside firms.
    \item Equilibrium adjustment determines who gains, who loses, and which effects persist.
\end{itemize}
\end{frame}

\begin{frame}{Four margins to separate}
\small
\begin{tabular}{p{0.22\linewidth} p{0.58\linewidth}}
\toprule
Margin & Labor-market meaning \\
\midrule
Invention & a new capability appears \\
Adoption & a firm, establishment, or worker begins using it \\
Diffusion & adoption spreads across firms, places, occupations, or workers \\
Reorganization & tasks, jobs, supervision, hiring, or workflow change around it \\
\bottomrule
\end{tabular}
\vspace{0.35cm}
\begin{itemize}
    \item A patent, robot installation, software rollout, and vacancy skill request are different events.
\end{itemize}
\end{frame}

\begin{frame}{Core framework}
\small
\[
\begin{aligned}
\text{Technology / Innovation}
&\rightarrow
\text{Task content}
\rightarrow
\text{Skill demand} \\
&\rightarrow
\text{Worker and firm adjustment}
\rightarrow
\text{Equilibrium outcomes}
\end{aligned}
\]
\begin{itemize}
    \item Tasks are the bridge from technology to labor demand.
    \item Skills matter because workers differ in comparative advantage across tasks.
    \item Firms decide whether to adopt, reorganize, hire, train, substitute, or expand.
    \item Workers decide whether to learn, move, switch jobs, bargain, or exit.
\end{itemize}
\end{frame}

\begin{frame}{Substitution, augmentation, and new tasks}
\small
\begin{tabular}{p{0.24\linewidth} p{0.37\linewidth} p{0.24\linewidth}}
\toprule
Channel & Direct task effect & Typical labor question \\
\midrule
Automation & tasks move from labor to capital or software & which workers are displaced? \\
Augmentation & workers become more productive in existing tasks & whose productivity rises? \\
New-task creation & labor gains new activities or occupations & where is labor reinstated? \\
\bottomrule
\end{tabular}
\vspace{0.35cm}
\begin{itemize}
    \item The same technology can operate through all three channels.
\end{itemize}
\end{frame}

\begin{frame}{Automation boundary}
\[
\frac{c_T}{a_{Tk}} \leq \frac{w}{a_{Lk}}
\]
\begin{itemize}
    \item $k$ indexes tasks.
    \item $c_T$ is the user cost of the technology.
    \item $a_{Tk}$ and $a_{Lk}$ are technology and labor productivity in task $k$.
    \item A shift in this boundary predicts task substitution, not the full equilibrium effect.
\end{itemize}
\end{frame}

\begin{frame}{Exposure is not effect}
\[
\Delta Y_{g t} =
\beta \, TechMeasure_{g t}
+ X_{g t}'\Gamma
+ \varepsilon_{g t}
\]
\begin{itemize}
    \item What does the technology measure capture: invention, exposure, adoption, diffusion, or realized use?
    \item What is $g$: task, occupation, worker, firm, establishment, industry, or place?
    \item What is $Y$: wages, employment, vacancies, skill demand, mobility, organization, or welfare?
    \item Which equilibrium margin is left outside the design?
\end{itemize}
\end{frame}

\begin{frame}{Measurement is a research frontier}
\small
\begin{tabular}{p{0.28\linewidth} p{0.34\linewidth} p{0.24\linewidth}}
\toprule
Measure & Closest object & Main warning \\
\midrule
Patents and patent text & invention and capability & not realized use \\
Robots & embodied automation adoption & narrow technology class \\
Software / IT & digital and organizational adoption & heterogeneous meaning \\
Vacancy text & desired task and skill demand & postings are not employment \\
Task surveys / O*NET & job content and exposure structure & can lag current work \\
Firm investment / resumes & realized adoption and skills & selection and access \\
\bottomrule
\end{tabular}
\end{frame}

\begin{frame}{Different measures, different objects}
\begin{itemize}
    \item Patent-text exposure can rank occupations by capability similarity before adoption is visible.
    \item Robot exposure can isolate one concrete embodied automation margin.
    \item Vacancy text can show changing employer demand for technology-related skills.
    \item Firm adoption data can move closer to realized use but make selection central.
    \item Disagreement across measures is often information, not a nuisance.
\end{itemize}
\end{frame}

\begin{frame}{Main literature map}
\small
\begin{tabular}{p{0.30\linewidth} p{0.46\linewidth}}
\toprule
Literature & Week 1 role \\
\midrule
Task content and SBTC & tasks connect technology to labor demand \\
Automation and new tasks & displacement, productivity, and reinstatement \\
Robots and ICT & concrete adoption margins \\
Patent/task exposure & capability-to-task measurement \\
Vacancies and firm adoption & desired skill demand and realized implementation \\
\bottomrule
\end{tabular}
\vspace{0.3cm}
\begin{itemize}
    \item The organizing question is always the same: what is measured and what labor outcome follows?
\end{itemize}
\end{frame}

\begin{frame}{Research design checklist}
\begin{enumerate}
    \item Name the technology margin.
    \item State whether the measure captures invention, exposure, adoption, diffusion, or realized use.
    \item Choose the unit: task, occupation, worker, firm, industry, or place.
    \item Name the labor outcome and adjustment horizon.
    \item State the equilibrium margin that remains offstage.
\end{enumerate}
\end{frame}

\begin{frame}{Week 1 lab}
\begin{itemize}
    \item Reproduce: construct a toy technology-to-task exposure score in the spirit of Webb.
    \item Diagnose: compare patent-text-style exposure, robot-style exposure, and vacancy-style skill demand.
    \item Transfer: move the logic to synthetic firms with different workforce task portfolios.
    \item Goal: understand how measurement choices create different empirical conclusions.
\end{itemize}
\end{frame}

\begin{frame}{Bridge to Week 2}
\begin{itemize}
    \item Week 1 builds the common language: technology, tasks, skills, adjustment, equilibrium, and measurement.
    \item Week 2 studies automation, AI, and labor demand more directly.
    \item The key next question: which tasks are displaced, complemented, or newly created when technologies diffuse?
\end{itemize}
\end{frame}

\begin{frame}{Takeaways}
\begin{itemize}
    \item Technology is a labor object when it changes tasks, skills, firms, wages, employment, or welfare.
    \item Invention, adoption, diffusion, and reorganization are separate empirical margins.
    \item Automation, augmentation, and new-task creation have different labor-market implications.
    \item Measurement is not a technical afterthought; it is one of the central open questions in the field.
\end{itemize}
\end{frame}

\end{document}
